\chapter{Torsion and the variation of mass}
\label{chap:mass}

For smooth K\"ahler potentials, Stokes' theorem makes the total Monge--Amp\`ere mass independent of the potential. With a nonclosed background, the same calculation leaves derivatives of powers of that form. The resulting mass polynomial identifies when invariance holds and gives explicit examples where it fails. This failure also explains why the global comparison principle must change in the Hermitian setting.

\section{Integration by parts with a non-closed form}

Let $X$ be a connected compact complex manifold of dimension $n\geq1$. We first compute the terms that remain when integration by parts differentiates a power of the background form.

\begin{lemma}\label{lem:ibp}
For a smooth real function $f$ and a smooth form $T$ of bidegree $(n-1,n-1)$,
\[
 \int_X\ddc f\wedge T=\int_X f\ddc T.
\]
\end{lemma}
\begin{proof}
Stokes' theorem applied first to $\mathrm{d}(\dc f\wedge T)$ gives
\[
 \int_X\ddc f\wedge T=\int_X\dc f\wedge \mathrm{d}T.
\]
The operator $\dc$ also has zero integral on forms of total degree $2n-1$: on the possible bidegrees this follows from Stokes' theorem for $\partial$ and $\bar\partial$. Applying this to $\dc(f\,\mathrm{d}T)$ and using $\dc \mathrm{d}=-\mathrm{d}\dc$ yields
\[
 0=\int_X\dc f\wedge \mathrm{d}T-\int_X f\ddc T.
\]
Combining the two identities proves the claim.
\end{proof}

\Needspace{10\baselineskip}
\begin{lemma}\label{lem:power-torsion}
For a smooth real $(1,1)$-form $\omega$ and an integer $q\geq1$,
\[
 \ddc(\omega^q)
 =q\,\omega^{q-1}\wedge\ddc\omega
 +q(q-1)\,\omega^{q-2}\wedge \mathrm{d}\omega\wedge\dc\omega,
\]
where the second term is omitted when $q=1$.
\end{lemma}
\begin{proof}
For $q=1$, the asserted formula is the identity $\ddc\omega=\ddc\omega$. Suppose $q\geq2$. Since $\omega$ has even degree, $\dc(\omega^q)=q\,\omega^{q-1}\wedge\dc\omega$. Applying $\mathrm{d}$ and again using the even degree of $\omega^{q-1}$ gives
\[
 \mathrm{d}\dc(\omega^q)
 =q(q-1)\,\omega^{q-2}\wedge \mathrm{d}\omega\wedge\dc\omega
 +q\,\omega^{q-1}\wedge \mathrm{d}\dc\omega.
\]
\end{proof}

The term involving $\mathrm{d}\omega\wedge\dc\omega$ is one reason that the condition $\ddc\omega=0$ must not be confused with closedness. In higher dimension it need not imply $\ddc(\omega^q)=0$.

\section{The mass polynomial}

Fix a Hermitian form $\omega$ on $X$. For $u\in C^\infty(X,\mathbb R)$ write $\omega_u=\omega+\ddc u$. Expanding its top power and integrating by parts will express the mass variation in terms of derivatives of powers of $\omega$. Set
\[
 \HH_\omega=\{u\in C^\infty(X,\R):\omega+\ddc u>0\},
 \qquad \Mass_\omega(u)=\int_X(\omega+\ddc u)^n.
\]
The set $\HH_\omega$ is convex, contains all constants, and is open in the $C^2$ topology.

\begin{theorem}\label{thm:mass-expansion}
For every smooth real function $u$,
\begin{equation}\label{eq:mass-expansion}
 \Mass_\omega(u)=\int_X\omega^n+
 \sum_{k=1}^{n-1}\binom nk
 \int_X u(\ddc u)^{k-1}\wedge\ddc(\omega^{n-k}).
\end{equation}
The equality is an identity of integrals and does not require $u\in\HH_\omega$.
\end{theorem}
\begin{proof}
Expand $(\omega+\ddc u)^n$ by the binomial theorem. The term with $k=n$ integrates to zero, since
\[
 (\ddc u)^n=\mathrm{d}\bigl(\dc u\wedge(\ddc u)^{n-1}\bigr).
\]
For $1\leq k\leq n-1$, apply \cref{lem:ibp} with
$T=(\ddc u)^{k-1}\wedge\omega^{n-k}$. Both $\mathrm{d}$ and $\dc$ annihilate $\ddc u$, so
\[
 \ddc T=(\ddc u)^{k-1}\wedge\ddc(\omega^{n-k}).
\]
This gives each of the terms in \eqref{eq:mass-expansion}. When $n=1$, the sum is empty and the mass is constant.
\end{proof}

\Needspace{9\baselineskip}
\begin{corollary}\label{cor:constant-mass-sufficient}
If $\ddc(\omega^j)=0$ for $1\leq j\leq n-1$, then
\[
 \Mass_\omega(u)=\int_X\omega^n
\]
for every smooth $u$. In particular, this holds when $\mathrm{d}\omega=0$.
\end{corollary}
\begin{proof}
Every term in the sum in \eqref{eq:mass-expansion} vanishes. A closed real $(1,1)$-form is annihilated by both $\partial$ and $\bar\partial$, hence so are its powers.
\end{proof}

\Needspace{11\baselineskip}
\begin{proposition}\label{prop:mass-variation}
For smooth $u,v$,
\begin{equation}\label{eq:mass-first-variation}
 \left.\frac{\mathrm{d}}{\mathrm{d}t}\right|_{t=0}\Mass_\omega(u+tv)
 =n\int_X v\ddc(\omega_u^{n-1}).
\end{equation}
If $\Mass_\omega$ is constant on $\HH_\omega$, then $\ddc(\omega_u^{n-1})=0$ for every $u\in\HH_\omega$.
\end{proposition}
\begin{proof}
Differentiating the finite polynomial gives
\[
 n\int_X\ddc v\wedge\omega_u^{n-1}.
\]
Apply \cref{lem:ibp}. For the second assertion, every smooth direction $v$ gives a line segment in $\HH_\omega$ for sufficiently small $|t|$, so the right side vanishes for all test functions $v$. The smooth top-degree form $\ddc(\omega_u^{n-1})$ must therefore vanish.
\end{proof}

\begin{corollary}\label{cor:surface-mass}
If $n=2$, then
\[
 \Mass_\omega(u)=\int_X\omega^2+2\int_Xu\ddc\omega.
\]
Consequently, $\Mass_\omega$ is constant on $\HH_\omega$ if and only if $\ddc\omega=0$.
\end{corollary}
\begin{proof}
The formula is \eqref{eq:mass-expansion} for $n=2$. It proves one implication directly; the other follows from \cref{prop:mass-variation} at $u=0$.
\end{proof}

A Hermitian form with $\ddc(\omega^{n-1})=0$ is called a Gauduchon form. Thus on a surface the Gauduchon condition is exactly the condition for smooth mass invariance. In dimension three, the expansion is
\[
 \Mass_\omega(u)=\int_X\omega^3
 +3\int_Xu\ddc(\omega^2)
 +3\int_Xu\ddc u\wedge\ddc\omega.
\]
The Gauduchon condition removes the middle term, but does not by itself remove the last one. Equivalently, \eqref{eq:mass-first-variation} vanishes at $u=0$, while mass invariance requires it to vanish at every admissible $u$. The full characterization of mass invariance is given in \cref{thm:mass-invariance-characterization}.

\Needspace{23\baselineskip}
\section{Mass variation on a complex torus}

The surface formula suggests choosing a background whose complex Hessian is nonzero and a potential that pairs nontrivially with it. A product of elliptic curves gives an explicit example in which the mass varies along a family of smooth potentials.

\begin{example}\label{ex:torus-mass}
Let $X=(\C/(\mathbb Z+\mathrm i\mathbb Z))^2$. For $j=1,2$, let $x_j,y_j$ be the real coordinates on its universal cover and put
\[
 z_j=x_j+\mathrm i y_j,
 \qquad \beta_j=\frac{\mathrm i}{2\pi}\,\mathrm{d}z_j\wedge \mathrm{d}\bar z_j.
\]
The forms $\beta_j$ descend to $X$. For a real parameter $t$, define
\[
 b(x_1)=2+\cos(2\pi x_1),\qquad
 \omega=\beta_1+b(x_1)\beta_2,\qquad
 u_t=t\cos(2\pi x_1).
\]
Since $1\leq b\leq3$, the form $\omega$ is Hermitian. It is not closed: $\mathrm{d}\omega=\mathrm{d}b\wedge\beta_2$ is not identically zero.

For a function $f$ of $x_1$ alone, with derivatives taken in $x_1$, one has $\ddc f=(f''/4)\beta_1$. Hence
\[
 \omega_{u_t}
 =\bigl(1-\pi^2t\cos(2\pi x_1)\bigr)\beta_1+b\beta_2>0
 \quad\text{if}\quad |t|<\pi^{-2}.
\]
Writing $B=\int_X\beta_1\wedge\beta_2>0$, direct expansion and integration give
\begin{equation}\label{eq:torus-mass}
 \Mass_\omega(u_t)=(4-\pi^2t)B.
\end{equation}
Indeed, the average of $\cos(2\pi x_1)$ is zero and that of its square is $1/2$, while
\[
 \omega_{u_t}^2
 =2\bigl(1-\pi^2t\cos(2\pi x_1)\bigr)
       \bigl(2+\cos(2\pi x_1)\bigr)\beta_1\wedge\beta_2.
\]
Thus even smooth potentials in the same singularity class have different masses.
\end{example}

The complex torus admits K\"ahler metrics, but the chosen equation form is nonclosed. Mass variation already occurs in this setting. The Hopf example in \cref{chap:hopf} will exhibit the further failure of an upper mass bound on a non-K\"ahler manifold.

\begin{corollary}\label{cor:mass-not-singularity-invariant}
Let $[u]$ denote the ordinary singularity type of $u$. In general, the rule
\[
 [u]\longmapsto\int_X\omega_u^n
\]
does not define a function of the ordinary singularity type, even if restricted to smooth strictly $\omega$-plurisubharmonic potentials. Nor is this mass determined by the multiplier ideals of the potential.
\end{corollary}
\begin{proof}
All the potentials $u_t$ in \cref{ex:torus-mass} are smooth on compact $X$, so they have the same singularity type. Their multiplier ideals are the structure sheaf for every positive multiple of the weight. Their masses differ by \eqref{eq:torus-mass}.
\end{proof}

Upper and lower masses taken over all representatives of an ordinary type are invariants of that type. Envelope constructions provide another way to choose a representative. The Gauduchon trace measures of \cref{thm:gaud-trace-homeomorphism} retain information about individual potentials. Their total mass, in contrast, depends only on the fixed background, by \cref{prop:gauduchon-trace}.

\Needspace{15\baselineskip}
\section{Comparison and mass invariance}

Adding a constant to a potential leaves its measure unchanged. Consequently, a comparison inequality on strict sublevel sets would compare total masses after a suitable constant shift. The varying-mass example therefore obstructs the usual global inequality.

\begin{proposition}\label{prop:comparison-forces-mass}
Suppose that a Hermitian form $\omega$ has the following property for every pair $u,v\in\HH_\omega$:
\begin{equation}\label{eq:naive-comparison}
 \int_{\{u<v\}}\omega_v^n
 \leq\int_{\{u<v\}}\omega_u^n.
\end{equation}
Then $\Mass_\omega$ is constant on $\HH_\omega$.
\end{proposition}
\begin{proof}
Fix $u,v\in\HH_\omega$. By compactness choose a constant $C$ such that $u<v+C$ everywhere. Since $\omega_{v+C}=\omega_v$, applying \eqref{eq:naive-comparison} to $u$ and $v+C$ gives $\Mass_\omega(v)\leq\Mass_\omega(u)$. Exchange $u$ and $v$ and repeat the argument to obtain the opposite inequality.
\end{proof}

\begin{corollary}\label{cor:comparison-fails}
The inequality \eqref{eq:naive-comparison} is false for general compact Hermitian backgrounds.
\end{corollary}
\begin{proof}
Use the torus $X$, form $\omega$, potentials $u_t$ and constant $B=\int_X\beta_1\wedge\beta_2$ of \cref{ex:torus-mass}. For $-\pi^{-2}<t<0$, set $u=0$ and $v=u_t+C$ with $C>|t|$. Then $\{u<v\}=X$, but
\[
 \int_X\omega_v^2=(4-\pi^2t)B>4B=\int_X\omega_u^2.
\]
\end{proof}

The modified comparison theorem in \cref{thm:modified-comparison} accounts for the derivatives of the background. Under the Guan--Li condition, the ordinary comparison principle is recovered.

\Needspace{30\baselineskip}
\section{Uniform mass bounds}

Once mass invariance fails, one can still ask for a uniform upper bound. Unlike invariance, finiteness of this bound is independent of the choice of Hermitian form. Let $\eta$ be a second Hermitian form on $X$.

Define the upper smooth mass
\[
 \Vsm(\omega)=\sup_{u\in\HH_\omega}\Mass_\omega(u)\in(0,+\infty].
\]
\begin{proposition}\label{prop:smooth-bmp-metric}
If $C>0$ is a real constant and $\omega\leq C\eta$, then
\[
 \Vsm(\omega)\leq C^n\Vsm(\eta).
\]
In particular, finiteness of $\Vsm$ is independent of the Hermitian form.
\end{proposition}
\begin{proof}
For $u\in\HH_\omega$, the function $u/C$ belongs to $\HH_\eta$, because $C\eta+\ddc u$ is positive and dominates $\omega+\ddc u$. For positive Hermitian matrices $A\leq B$, all eigenvalues of $A^{-1/2}BA^{-1/2}$ are at least one, so $\det A\leq\det B$. Applying this pointwise gives monotonicity of the top exterior power. Therefore
\[
 \int_X(\omega+\ddc u)^n
 \leq C^n\int_X\bigl(\eta+\ddc(u/C)\bigr)^n
 \leq C^n\Vsm(\eta).
\]
Take the supremum over $u$. Any two Hermitian forms on compact $X$ are comparable in both directions, proving the last assertion.
\end{proof}
